% Options for packages loaded elsewhere
\PassOptionsToPackage{unicode}{hyperref}
\PassOptionsToPackage{hyphens}{url}
\documentclass[
]{article}
\usepackage{xcolor}
\usepackage{amsmath,amssymb}
\setcounter{secnumdepth}{-\maxdimen} % remove section numbering
\usepackage{iftex}
\ifPDFTeX
  \usepackage[T1]{fontenc}
  \usepackage[utf8]{inputenc}
  \usepackage{textcomp} % provide euro and other symbols
\else % if luatex or xetex
  \usepackage{unicode-math} % this also loads fontspec
  \defaultfontfeatures{Scale=MatchLowercase}
  \defaultfontfeatures[\rmfamily]{Ligatures=TeX,Scale=1}
\fi
\usepackage{lmodern}
\ifPDFTeX\else
  % xetex/luatex font selection
\fi
% Use upquote if available, for straight quotes in verbatim environments
\IfFileExists{upquote.sty}{\usepackage{upquote}}{}
\IfFileExists{microtype.sty}{% use microtype if available
  \usepackage[]{microtype}
  \UseMicrotypeSet[protrusion]{basicmath} % disable protrusion for tt fonts
}{}
\makeatletter
\@ifundefined{KOMAClassName}{% if non-KOMA class
  \IfFileExists{parskip.sty}{%
    \usepackage{parskip}
  }{% else
    \setlength{\parindent}{0pt}
    \setlength{\parskip}{6pt plus 2pt minus 1pt}}
}{% if KOMA class
  \KOMAoptions{parskip=half}}
\makeatother
\setlength{\emergencystretch}{3em} % prevent overfull lines
\providecommand{\tightlist}{%
  \setlength{\itemsep}{0pt}\setlength{\parskip}{0pt}}
\usepackage{bookmark}
\IfFileExists{xurl.sty}{\usepackage{xurl}}{} % add URL line breaks if available
\urlstyle{same}
\hypersetup{
  hidelinks,
  pdfcreator={LaTeX via pandoc}}

\author{}
\date{}

\begin{document}

\section{\texorpdfstring{Coordinate-Wise Additivity and the \(\ell^1\)
Norm on Finite Graph
Cochains}{Coordinate-Wise Additivity and the \textbackslash ell\^{}1 Norm on Finite Graph Cochains}}\label{coordinate-wise-additivity-and-the-ell1-norm-on-finite-graph-cochains}

\textbf{Author:} Jeremy H. Carroll \textbf{Date:} March 2026
\textbf{Version:} 0.2 (Pre-print)

\begin{center}\rule{0.5\linewidth}{0.5pt}\end{center}

\subsection{Abstract}\label{abstract}

We prove that on finite directed graphs, any system of seminorms on the
cochain space
\(C^1(G, \mathbb{R}^k) = \bigoplus_{e \in E} \mathbb{R}^k\) satisfying
edge locality, coordinate additivity over disjoint supports, coordinate
symmetry, and graph invariance must be proportional to the \(\ell^1\)
norm, up to a global scalar. This characterizes \(\ell^1\) as the unique
geometry compatible with separable defect aggregation on finite
cochains, providing a combinatorial foundation for the
functional-analytic and dynamical classifications established in
companion work.

\textbf{Conventions.} Norms are denoted \(\lVert \cdot \rVert\) with
appropriate subscripts; scalar absolute values are denoted \(|\cdot|\).

\begin{center}\rule{0.5\linewidth}{0.5pt}\end{center}

\subsection{1. Introduction}\label{introduction}

This paper is the first step in a sequence analyzing projection-induced
defects on presheaf coboundaries through progressively richer structural
settings:

\begin{itemize}
\tightlist
\item
  \textbf{Paper 002 (this paper)} establishes that coordinate-wise
  additivity and symmetry on finite graph cochains force an \(\ell^1\)
  geometry. The result is purely combinatorial and requires no
  functional-analytic or dynamical structure.
\item
  \textbf{Paper 001} {[}7{]} extends this classification to Banach
  presheaves, showing that functoriality under bounded linear maps
  similarly forces the \(\ell^1\) coboundary norm as the unique additive
  diagnostic compatible with the ambient Banach geometry.
\item
  \textbf{Paper 000} {[}6{]} applies this rigidity to dynamics, proving
  that projection of linear evolution onto a nonlinear manifold produces
  intrinsic coboundary defects whose natural measurement is therefore
  necessarily \(\ell^1\).
\end{itemize}

The purpose of the present paper is to isolate the minimal structural
mechanism behind this rigidity. Rather than beginning with Banach space
assumptions or dynamical systems, we show that \(\ell^1\) geometry
already emerges at the level of finite graph cochains once one requires
separability across independent coordinates and edges. This establishes
the combinatorial core of the classification, which is then lifted and
applied in subsequent work.

A central mathematical question in obstruction theories is why the local
diagnostic norm itself, operating on the \(k\)-dimensional measurement
channels (the stalks), should assume an \(\ell^1\) geometry instead of
\(\ell^2\) or \(\ell^\infty\). While it is a classical result in Banach
lattice theory that norms additive over disjoint supports correspond to
\(\ell^1\)-type structures (e.g., Kakutani's abstract \(L\)-spaces
{[}1{]}), this paper explicitly localizes that principle to finite graph
cochains to prove global geometric rigidity. When forced to respect both
axes of separability --- edges and coordinates --- the diagnosis becomes
uniquely forced.

\begin{center}\rule{0.5\linewidth}{0.5pt}\end{center}

\subsection{\texorpdfstring{2. \(\ell^1\) Characterization on Graph
Cochains}{2. \textbackslash ell\^{}1 Characterization on Graph Cochains}}\label{ell1-characterization-on-graph-cochains}

We recall that a \textbf{seminorm} on a vector space \(V\) is a
functional \(\nu: V \to \mathbb{R}_{\ge0}\) that is subadditive
(\(\nu(x+y) \le \nu(x) + \nu(y)\)) and absolutely homogeneous
(\(\nu(cx) = |c|\nu(x)\) for all scalars \(c \in \mathbb{R}\)).

\textbf{Theorem 2.1.} For any finite directed graph \(G=(V,E)\),
consider the cochain space
\[ C^1(G,\mathbb{R}^k)=\bigoplus_{e\in E}\mathbb{R}^k. \]

Let \(\nu_G:C^1(G,\mathbb{R}^k)\to\mathbb{R}_{\ge0}\) be a system of
seminorms, defined uniformly across all finite graphs, satisfying:

\begin{enumerate}
\def\labelenumi{\arabic{enumi}.}
\tightlist
\item
  \textbf{Edge locality (local additivity)}
  \[ \nu_G(\delta)=\sum_{e\in E}\nu_e(\delta_e) \]
\item
  \textbf{Coordinate symmetry}
  \[ \nu_e(x_{\sigma(1)},\dots,x_{\sigma(k)}) = \nu_e(x_1,\dots,x_k) \]
  for every permutation \(\sigma\).
\item
  \textbf{Coordinate additivity (local additivity over disjoint
  supports)} If \(x,y\in\mathbb{R}^k\) have disjoint coordinate support
  (\(\mathrm{supp}(x) \cap \mathrm{supp}(y) = \varnothing\) where
  \(\mathrm{supp}(x) = \{i : x_i \neq 0\}\)), then
  \[ \nu_e(x+y)=\nu_e(x)+\nu_e(y). \]
\item
  \textbf{Graph invariance} For any graph isomorphism
  \(\sigma:G\to G'\), \[ \nu_{G'}(\sigma_*\delta)=\nu_G(\delta). \]
\end{enumerate}

\emph{Remark.} Axiom 3 is precisely the defining structural property of
abstract \(L\)-spaces: additivity over disjoint supports {[}1{]}.

Then there exists \(\alpha > 0\) such that universally for all \(G\):
\[ \nu_G(\delta)=\alpha\sum_{e\in E}\|\delta_e\|_1. \]

\textbf{Proof.}

By edge locality it suffices to classify the seminorm \(\nu_e\) on
\(\mathbb{R}^k\).

Let \(x=(x_1,\dots,x_k)\). Decompose \[ x=\sum_{i=1}^k x_i \hat{e}_i. \]

The vectors \(x_i \hat{e}_i\) have disjoint coordinate support, so
coordinate additivity gives
\[ \nu_e(x)=\sum_{i=1}^k\nu_e(x_i \hat{e}_i). \]

Since \(\nu_e\) is a seminorm, it is absolutely homogeneous, providing
\[ \nu_e(x_i \hat{e}_i)=|x_i|\nu_e(\hat{e}_i). \]

By coordinate symmetry, \[ \nu_e(\hat{e}_i)=w_e \] for all \(i\).

Hence \[ \nu_e(x)=w_e\sum_{i=1}^k |x_i|=w_e\|x\|_1. \]

Thus across the entire graph,
\[ \nu_G(\delta)=\sum_{e\in E} w_e\|\delta_e\|_1. \]

Finally, graph invariance implies all intrinsic weights coincide.
Because the construction is required to be invariant across all finite
graphs, and there exist graphs whose automorphism groups act
transitively on edges (e.g., complete graphs or regular cycles), the
weights must coincide universally. Specifically, for any two edges
\(e_1, e_2\) (possibly in different graphs), embed both into a common
graph \(G'\) where a graph automorphism maps \(e_1\) to \(e_2\); by
invariance, \(w_{e_1} = w_{e_2}\). Thus: \[ w_e=\alpha. \]

Therefore \[ \nu_G(\delta)=\alpha\sum_{e\in E}\|\delta_e\|_1. \]
\(\square\)

\textbf{Corollary 2.2 (Discrete \(L^1\) Structure of the Defect Space).}
Under the hypotheses of Theorem 2.1, the cochain space equipped with the
diagnostic seminorm satisfies

\[
(C^1(G,\mathbb{R}^k),\nu_G)
\cong
L^1(E\times\{1,\dots,k\})
\]

with respect to the counting measure.

\textbf{Proof.}

By Theorem 2.1,

\[
\nu_G(\delta)
=
\alpha\sum_{e\in E}\|\delta_e\|_1.
\]

Writing \(\delta_{e,i}\) for the \(i\)-th coordinate of \(\delta_e\), we
obtain

\[
\nu_G(\delta)
=
\alpha\sum_{e\in E}\sum_{i=1}^k |\delta_{e,i}|.
\]

Define

\[
f:E\times\{1,\dots,k\}\to\mathbb{R}
\]

by

\[
f(e,i)=\delta_{e,i}.
\]

Then

\[
\nu_G(\delta)
=
\alpha\|f\|_{L^1}.
\]

Thus the diagnostic space is isometrically equivalent to
\(L^1(E\times\{1,\dots,k\})\) under the counting measure. \(\square\)

\begin{center}\rule{0.5\linewidth}{0.5pt}\end{center}

\subsection{3. Structural Consequences}\label{structural-consequences}

\subsubsection{\texorpdfstring{3.1 Coarea Duality
(\(k = 1\))}{3.1 Coarea Duality (k = 1)}}\label{coarea-duality-k-1}

For scalar signals, the classified norm has a distinguished topological
property.

\textbf{Fact 3.1 (Discrete Coarea Formula} {[}2, 3{]}\textbf{).} For
\(v: V \to \mathbb{R}\):
\[ \text{TV}_G(v) = \int_{-\infty}^{\infty} |\partial V_t| \, dt \]
where \(V_t = \{m : v_m > t\}\) and \(|\partial V_t|\) counts edges
crossing the boundary.

When boundary conditions force a binary partition
(\(v_a = 1, v_b = 0\)), \(\ell^1\) minimization recovers the minimum
\(s\)-\(t\) cut via max-flow/min-cut {[}5{]}.

\textbf{Remark 3.2.} For \(k > 1\), the total variation
\(\sum_e \|v_n - v_m\|_1\) decomposes (by Theorem 2.1) as a sum of \(k\)
independent scalar total variations, each individually satisfying the
coarea formula.

\subsubsection{3.2 The Geometry of the Defect
Field}\label{the-geometry-of-the-defect-field}

The explicit inclusion of \textbf{Coordinate Additivity} formally
captures the physical independence of disparate measurement channels.
Permutation symmetry and absolute homogeneity alone permit alternative
norms such as \(\ell^\infty\) or \(\ell^2\). The strict structural
requirement that physically isolated defect channels --- the orthogonal
coordinates --- must accumulate their structural signatures linearly
thereby uniquely determines the \(\ell^1\) stalk geometry under these
constraints.

Thus, the unique seminorm compatible with the stated structural
principles under the imposed local additivity and symmetry constraints
is the \(\ell^1\) coboundary norm: an \(\ell^1\) geometry over the
product index set.

Conceptually, our theorem establishes (cf.~{[}4{]}) that the defect
space restricts to:
\[ C^1(G,\mathbb{R}^k) \cong L^1(E \times \{1,\dots,k\}) \] whenever the
diagnostic respects independent components. This serves as a finite
combinatorial analogue of Kakutani's (1941) concrete characterization of
abstract \(L\)-spaces {[}1{]}. Local additivity across independent
structural components forces the \(\ell^1\) geometry over the counting
measure. Defect fields thus behave strictly as discrete 1-forms whose
absolute total variation tracks graph cuts and sparse obstructions
{[}3{]}.

\subsubsection{3.3 Role in the Broader Classification
Program}\label{role-in-the-broader-classification-program}

The result obtained here isolates the combinatorial origin of \(\ell^1\)
geometry as a consequence of coordinate-wise and edge-wise separability.
In subsequent work, this principle is shown to be stable under
significant generalization. In particular, when the cochain spaces are
replaced by Banach presheaves and coordinate additivity is replaced by
structural compatibility in the form of functoriality under bounded
linear maps {[}7{]}, the same \(\ell^1\) structure is recovered as the
unique compatible seminorm.

This stability across levels of abstraction indicates that the
appearance of \(\ell^1\) is not an artifact of discrete coordinates, but
a structural feature of additive defect aggregation. The present theorem
therefore provides the minimal case from which the more general Banach
{[}7{]} and dynamical {[}6{]} results can be understood as extensions.

\emph{Remark (Observer-theoretic grounding).} The axioms of Theorem 2.1
admit an operational interpretation: they are the minimal conditions
under which a local agent can consistently detect and aggregate defects
across a graph. Edge locality requires that defects be detectable from a
single edge; coordinate additivity requires that independent measurement
channels contribute additively without cancellation; symmetry requires
that the diagnostic cannot distinguish structurally identical
subsystems. Under this reading, the \(\ell^1\) norm is not merely a
mathematical consequence of abstract axioms, but the unique diagnostic
compatible with consistent local observation. This interpretation is
developed in the companion paper {[}6{]}.

\begin{center}\rule{0.5\linewidth}{0.5pt}\end{center}

\subsection{Acknowledgments}\label{acknowledgments}

No external funding. No conflicts of interest.

\begin{center}\rule{0.5\linewidth}{0.5pt}\end{center}

\subsection{References}\label{references}

\begin{enumerate}
\def\labelenumi{\arabic{enumi}.}
\tightlist
\item
  Kakutani, S. (1941). \emph{Concrete representation of abstract
  \((L)\)-spaces and the mean ergodic theorem}. Annals of Mathematics,
  42(2), 523-537.
\item
  Fleming, W. H., \& Rishel, R. (1960). \emph{An integral formula for
  total gradient variation}. Archiv der Mathematik, 11(1), 218-222.
\item
  Chambolle, A., et al.~(2010). \emph{An introduction to total variation
  for image analysis.} In: Theoretical Foundations and Numerical Methods
  for Sparse Recovery. De Gruyter.
\item
  Rockafellar, R. T. (1970). \emph{Convex Analysis.} Princeton
  University Press.
\item
  Ford, L. R., \& Fulkerson, D. R. (1956). \emph{Maximal flow through a
  network.} Canadian Journal of Mathematics, 8, 399-404.
\item
  Carroll, J. H. (2026). \emph{Projection Obstruction Theory: Retraction Nonlinearity, \(\ell^1\) Rigidity, and Density Scaling Rigidity, and Density Scaling.} Zenodo.
  https://doi.org/10.5281/zenodo.191151803
\item
  Carroll, J. H. (2026). \emph{The Free \(\ell^1\) Seminorm on Banach
  Presheaf Coboundaries.} Zenodo.
  https://doi.org/10.5281/zenodo.191152115
\end{enumerate}

\end{document}
